\documentclass[12pt]{article}
\usepackage[utf8]{inputenc}
\usepackage[T1]{fontenc}
\usepackage[margin=1in]{geometry}
\usepackage{setspace}
\usepackage{xcolor}
\usepackage[colorlinks=true, urlcolor=blue!60!black, linkcolor=blue!60!black]{hyperref}
\setlength{\parindent}{2em}
\setlength{\parskip}{1.25em}
\renewcommand{\baselinestretch}{1.0}
\begin{document}
\begin{center}
{\LARGE \textit{DeShaney v. Winnebago County Department of Social Services}, 489 U.S. 189 (1989)\par}
\vspace{0.8em}
{\large Case Brief\par}
\end{center}
\vspace{1em}

\section{Memory Jogger}
A Wisconsin county social-services agency repeatedly documented signs that four-year-old Joshua DeShaney was being abused by his father but never removed him; the Court held 6--3 that the Due Process Clause imposes no duty to protect against private violence unless the State has taken the person into custody or created the danger (\href{https://supreme.justia.com/cases/federal/us/489/189/}{Justia}).

\section{Detailed Case Facts}
Joshua DeShaney was born in 1979. Following his parents' 1980 divorce in Wyoming, custody was awarded to his father, Randy DeShaney, who moved with Joshua to Neenah, in Winnebago County, Wisconsin (\href{https://supreme.justia.com/cases/federal/us/489/189/}{489 U.S. at 191}). In January 1982, Randy's second wife reported to police during their divorce that Randy had ``hit the boy, causing marks, and [was] a prime case for child abuse''; the Winnebago County Department of Social Services (DSS) interviewed Randy, accepted his denial, and took no further action (\href{https://supreme.justia.com/cases/federal/us/489/189/}{489 U.S. at 192}).

In January 1983, Joshua was hospitalized with multiple bruises and abrasions. Suspecting abuse, the examining physician notified DSS, which obtained a juvenile court order placing Joshua in the hospital's temporary custody (\href{https://supreme.justia.com/cases/federal/us/489/189/}{489 U.S. at 192}). Three days later, a county ``Child Protection Team'' concluded there was insufficient evidence to retain custody and recommended returning Joshua to Randy under a voluntary agreement involving preschool enrollment, counseling, and having Randy's girlfriend move out. The juvenile court dismissed the case and returned Joshua home (\href{https://supreme.justia.com/cases/federal/us/489/189/}{489 U.S. at 192}).

Over the next year, emergency room staff twice reported new suspicious injuries to DSS (in February and November 1983), and during monthly home visits a DSS caseworker observed recurring head injuries, noted Randy's noncompliance with the voluntary agreement, and recorded continuing suspicions of abuse in her files---yet DSS took no action, even when the caseworker was twice told Joshua was too ill to see her (\href{https://supreme.justia.com/cases/federal/us/489/189/}{489 U.S. at 192--193}). In March 1984, Randy beat four-year-old Joshua so severely that he fell into a coma and suffered catastrophic, permanent brain damage, leaving him institutionalized for life; Randy was convicted of child abuse (\href{https://supreme.justia.com/cases/federal/us/489/189/}{489 U.S. at 193}).

\section{Procedural History}
Joshua and his mother sued Winnebago County, DSS, and individual DSS employees under 42 U.S.C. \S~1983 in the U.S. District Court for the Eastern District of Wisconsin, alleging that their failure to intervene violated Joshua's Fourteenth Amendment substantive due process liberty interest in bodily integrity (\href{https://supreme.justia.com/cases/federal/us/489/189/}{489 U.S. at 193}). The District Court granted summary judgment for respondents, the Seventh Circuit affirmed (812 F.2d 298 (1987)), and the Supreme Court granted certiorari to resolve a circuit split (\href{https://supreme.justia.com/cases/federal/us/489/189/}{489 U.S. at 193--194}).

\section{Judicial Votes}
The Supreme Court \textbf{affirmed} the Seventh Circuit by a \textbf{6--3} vote. The case was argued November 2, 1988, and decided February 22, 1989 (\href{https://supreme.justia.com/cases/federal/us/489/189/}{Justia}; \href{https://www.law.cornell.edu/supremecourt/text/489/189}{Cornell LII}):

\begin{itemize}
    \item \textbf{Majority Opinion:} Chief Justice Rehnquist delivered the opinion of the Court, joined by Justices White, Stevens, O'Connor, Scalia, and Kennedy.
    \item \textbf{Concurring Opinions:} None.
    \item \textbf{Dissenting Opinions:} Justice Brennan filed a dissenting opinion, joined by Justices Marshall and Blackmun. Justice Blackmun filed a dissenting opinion.
\end{itemize}

\section{Holding}
\textbf{Question Presented:}

Does a state or local social services agency violate the substantive component of the Fourteenth Amendment's Due Process Clause when it knows a child faces a serious risk of abuse by a private custodial parent, takes steps to monitor the child, and fails to intervene to prevent severe private violence?

\textbf{Holding:} \textbf{No}; judgment \textbf{affirmed} (6--3) (\href{https://supreme.justia.com/cases/federal/us/489/189/}{489 U.S. at 191, 203}). The Due Process Clause is a charter of negative liberties---a limitation on the State's power to act, not a guarantee of minimal levels of safety and security. It confers no affirmative right to governmental aid or protection against private violence unless the State has restrained an individual's liberty through custody (creating a ``special relationship'') or affirmatively created or worsened the danger (\href{https://supreme.justia.com/cases/federal/us/489/189/}{489 U.S. at 195--201}).

\section{Analysis of Opinions}
\subsection{Majority Opinion (Chief Justice Rehnquist)}

\subsubsection{Text, History, and the Negative-Liberties Baseline}

Chief Justice Rehnquist began with the text of the Fourteenth Amendment: ``[n]o State shall . . . deprive any person of life, liberty, or property, without due process of law'' (\href{https://supreme.justia.com/cases/federal/us/489/189/}{489 U.S. at 194--195}). Because petitioners claimed a categorical substantive entitlement to protection rather than procedural safeguards (and waived any statutory procedural-entitlement claim below, \href{https://supreme.justia.com/cases/federal/us/489/189/}{489 U.S. at 195 n.2}), the Court held that neither text nor history supports an affirmative governmental duty to prevent private harm. The Clause was designed ``to protect the people from the State, not to ensure that the State protected them from each other,'' leaving the scope of protective services to democratic politics (\href{https://supreme.justia.com/cases/federal/us/489/189/}{489 U.S. at 196}). Consistent with \textit{Harris v. McRae}, 448 U.S. 297 (1980), and \textit{Lindsey v. Normet}, 405 U.S. 56 (1972), because the State has no constitutional duty to provide protective services, it cannot be liable under the Due Process Clause for injuries those services might have averted---absent selective denial violating Equal Protection (\href{https://supreme.justia.com/cases/federal/us/489/189/}{489 U.S. at 196--197 \& n.3}).

\subsubsection{Narrowing the Custodial ``Special Relationship'' Exception}

Petitioners argued that DSS's awareness of Joshua's peril and voluntary assumption of monitoring created a ``special relationship'' imposing an affirmative duty of care (\href{https://supreme.justia.com/cases/federal/us/489/189/}{489 U.S. at 197}). The Court rejected that reading of \href{https://supreme.justia.com/cases/federal/us/429/97/}{\textit{Estelle v. Gamble}, 429 U.S. 97 (1976)} (medical care for prisoners), and \href{https://supreme.justia.com/cases/federal/us/457/307/}{\textit{Youngberg v. Romeo}, 457 U.S. 307 (1982)} (reasonable safety for involuntarily committed mental patients). Those cases establish only that when the State holds a person in custody against his will---through incarceration, institutionalization, or similar restraint---rendering him unable to care for himself, the State's affirmative restraint triggers a duty to meet basic human needs (\href{https://supreme.justia.com/cases/federal/us/489/189/}{489 U.S. at 198--200}).

Here, Joshua was injured while in the custody of his natural father, a private actor, not in state custody (leaving open whether state-operated foster care is analogous to custody, \href{https://supreme.justia.com/cases/federal/us/489/189/}{489 U.S. at 201 \& n.9}). Although Wisconsin briefly held temporary custody at the hospital, returning Joshua to his father ``placed him in no worse position than that in which he would have been had it not acted at all'' (\href{https://supreme.justia.com/cases/federal/us/489/189/}{489 U.S. at 201}). Crucially, the State ``played no part in [the dangers'] creation, nor did it do anything to render him any more vulnerable to them'' (\href{https://supreme.justia.com/cases/federal/us/489/189/}{489 U.S. at 201}).

\subsubsection{State Tort Law and the Social-Worker Dilemma}

Rehnquist emphasized that the Fourteenth Amendment ``does not transform every tort committed by a state actor into a constitutional violation'' (\href{https://supreme.justia.com/cases/federal/us/489/189/}{489 U.S. at 202}). Wisconsin remained free to impose tort liability through its courts or legislature. Moreover, had DSS removed Joshua prematurely, it would have faced Due Process claims for intruding into the parent-child relationship (\href{https://supreme.justia.com/cases/federal/us/489/189/}{489 U.S. at 202--203}).

\subsection{Dissenting Opinion (Justice Brennan)}

Justice Brennan criticized the majority's rigid baseline distinguishing action from inaction. Focusing on Wisconsin's affirmative acts, he argued that Wisconsin's statutory child-welfare regime funneled all abuse reports from police, doctors, and neighbors exclusively to DSS, effectively displacing private and other public sources of rescue. By monopolizing protection and then failing to act despite eerie, documented knowledge (the caseworker later remarked, ``I just knew the phone would ring some day and Joshua would be dead''), DSS actively confined Joshua within a violent home (\href{https://supreme.justia.com/cases/federal/us/489/189/}{489 U.S. at 203--212}).

\subsection{Dissenting Opinion (Justice Blackmun)}

Writing separately, Justice Blackmun condemned the majority's ``sterile formalism,'' analogizing its claim of doctrinal compulsion to antebellum judges enforcing the Fugitive Slave Acts. Urging a ``sympathetic'' reading of the Fourteenth Amendment grounded in fundamental justice, he famously lamented: ``Poor Joshua!'' (\href{https://supreme.justia.com/cases/federal/us/489/189/}{489 U.S. at 212--213}).

\subsection{Subsequent Developments}

\begin{itemize}

\item \textbf{Reinforced and Extended (\textit{Collins}; \textit{Castle Rock}):} The Supreme Court reinforced this boundary in \href{https://www.law.cornell.edu/supremecourt/text/503/115}{\textit{Collins v. City of Harker Heights}, 503 U.S. 115 (1992)} (no substantive due process duty to provide municipal employees a safe workplace), and closed the procedural-entitlement loophole left open in \textit{DeShaney}'s footnote 2 in \href{https://supreme.justia.com/cases/federal/us/545/748/}{\textit{Town of Castle Rock v. Gonzales}, 545 U.S. 748 (2005)}, holding that even a seemingly mandatory state domestic-violence restraining-order statute confers no constitutionally protected ``property'' interest in police enforcement.

\item \textbf{State-Created Danger in the Circuits:} Meanwhile, Chief Justice Rehnquist's observation that Wisconsin ``played no part in [the dangers'] creation, nor did it do anything to render him any more vulnerable to them'' (\href{https://supreme.justia.com/cases/federal/us/489/189/}{489 U.S. at 201}) spawned the federal circuits' \textbf{state-created danger} doctrine. Most federal courts of appeals recognize substantive due process liability when a state actor affirmatively creates or enhances a foreseeable risk of private harm with conscience-shocking culpability, though the Supreme Court has never formally endorsed the doctrine and circuit tests diverge sharply (\href{https://lawreview.uchicago.edu/online-archive/exploring-statutory-remedies-state-created-dangers-after-fisher-v-moore}{U. Chi. L. Rev. Online}).

\end{itemize}

\section{Examples: Future Applications}
\begin{enumerate}

\item \textbf{Same-Side (State Prevails): Slow Response to Protective-Order Violations.} City police receive three 911 calls reporting that a local resident is threatening his estranged spouse outside her apartment in violation of a state protective order. Dispatch logs the calls as low priority, no patrol car arrives for two hours, and the resident assaults the spouse. Under \textit{DeShaney} and \href{https://supreme.justia.com/cases/federal/us/545/748/}{\textit{Castle Rock}}, a \S~1983 Due Process suit fails: the victim was not in state custody, police did not affirmatively create the danger, and state protective-order statutes confer no federal procedural property interest in arrest (remedies lie, if at all, in state tort law or Equal Protection if non-enforcement was discriminatory).

\item \textbf{Same-Side (State Prevails): Broken Promises of Extra Patrols.} A woman tells her local police department that a former boyfriend has threatened to kill her. An officer promises that patrol cars will drive past her house every hour that week, but none do, and on the third night the former boyfriend breaks in and injures her. She sues the officer and the city under \S~1983, arguing that the promise created a special relationship. Under \textit{DeShaney}, the claim fails. The State did not restrain her liberty, and the danger came from a private actor. The Court made clear that an affirmative duty arises not from the State's knowledge of a person's predicament or its ``expressions of intent to help him,'' but only from limits it has imposed on the person's freedom to act on her own behalf. Any remedy lies in state tort law (\href{https://supreme.justia.com/cases/federal/us/489/189/}{Justia}).

\item \textbf{Opposite-Side (Plaintiff Prevails): A Child Stranded on the Highway.} After arresting an intoxicated driver at 1:00 a.m. in sub-freezing weather on an unlit rural interstate, a state trooper confiscates the car keys and orders the driver's sober seven-year-old passenger out of the car to ``walk home'' along the highway shoulder rather than waiting for a guardian. The child is struck by a passing motorist. Unlike \textit{DeShaney}, the officer did not merely stand by while a preexisting private danger unfolded. By affirmatively removing the child from a safe vehicle and stranding her in an exposed, lethal environment she could not escape, the state actor affirmatively created the peril and rendered the child far more vulnerable to harm (\href{https://supreme.justia.com/cases/federal/us/489/189/}{489 U.S. at 201}; cf. \textit{White v. Rochford}, 592 F.2d 381 (7th Cir. 1979)). In circuits recognizing state-created danger (and arguably as a brief custodial deprivation of liberty), the plaintiff states a viable substantive due process claim.

\item \textbf{Opposite-Side (Custody Triggers a Duty): A Prisoner Beaten by a Cellmate.} A state prisoner tells prison officers that his cellmate has threatened to kill him and asks to be moved. The officers laugh off the request, and that night the cellmate beats him severely. Unlike Joshua, the prisoner is in state custody: the State has locked him in a cell with his attacker and taken away his ability to protect himself or seek help. \textit{DeShaney} itself explained that when the State so restrains a person's liberty, as in \textit{Estelle v.\ Gamble} and \textit{Youngberg v.\ Romeo}, it assumes a duty to provide for his basic needs, including reasonable safety. The prisoner can state a constitutional claim if he shows the officers knew of and disregarded a substantial risk to his safety, the Eighth Amendment standard the Court applied to prisoners attacked by other inmates in \textit{Farmer v.\ Brennan}, 511 U.S.\ 825 (1994).

\item \textbf{Fence-Sitter (Foster Care): Abuse in a State-Selected Foster Home.} A county child-welfare agency removes a six-year-old from her parents and places her with foster parents it has licensed and supervises. Caseworkers receive two reports that the foster father is hitting her but do not investigate, and she is seriously injured. The agency relies on \textit{DeShaney}: the abuser was a private person, not a state employee. The child answers that the State took her from her family and chose where she would live, a restraint on her liberty much closer to institutionalization than Joshua's situation. \textit{DeShaney} expressly left open whether a child in state-arranged foster care stands in a custodial relationship with the State (489 U.S.\ at 201 n.9), and the Supreme Court has not resolved the question, so the outcome is uncertain.

\end{enumerate}

\section{Critique}
\textbf{Action vs. Inaction Formalism:} Critics echo Justice Brennan's insight that ``inaction'' is a baseline choice: Wisconsin affirmatively enacted custody laws giving Randy exclusive legal control over Joshua and a mandatory reporting monopoly directing all rescue efforts to DSS.

\textbf{Reconstruction History Tension:} Placed in Chapter 5, \textit{DeShaney} stands in sharp tension with the Fourteenth Amendment's Reconstruction history, which was enacted largely because Southern states refused to protect Black citizens from private Ku Klux Klan violence---though Rehnquist addressed selective non-protection solely under the Equal Protection Clause (\href{https://supreme.justia.com/cases/federal/us/489/189/}{489 U.S. at 197 n.3}).

\section{Key Quotations}
\begin{quote}
``The Clause is phrased as a limitation on the State's power to act, not as a guarantee of certain minimal levels of safety and security. It forbids the State itself to deprive individuals of life, liberty, or property without `due process of law,' but its language cannot fairly be extended to impose an affirmative obligation on the State to ensure that those interests do not come to harm through other means. . . . Its purpose was to protect the people from the State, not to ensure that the State protected them from each other.''
\end{quote}
\noindent --- Chief Justice Rehnquist, for the Court, \href{https://supreme.justia.com/cases/federal/us/489/189/}{489 U.S. at 195--196}; \href{https://www.law.cornell.edu/supremecourt/text/489/189}{Cornell LII}.

\begin{quote}
``The affirmative duty to protect arises not from the State's knowledge of the individual's predicament or from its expressions of intent to help him, but from the limitation which it has imposed on his freedom to act on his own behalf.''
\end{quote}
\noindent --- Chief Justice Rehnquist, for the Court, \href{https://supreme.justia.com/cases/federal/us/489/189/}{489 U.S. at 200}; \href{https://www.law.cornell.edu/supremecourt/text/489/189}{Cornell LII}.

\begin{quote}
``While the State may have been aware of the dangers that Joshua faced in the free world, it played no part in their creation, nor did it do anything to render him any more vulnerable to them.''
\end{quote}
\noindent --- Chief Justice Rehnquist, for the Court, \href{https://supreme.justia.com/cases/federal/us/489/189/}{489 U.S. at 201}; \href{https://www.law.cornell.edu/supremecourt/text/489/189}{Cornell LII}.

\begin{quote}
``Poor Joshua! Victim of repeated attacks by an irresponsible, bullying, cowardly, and intemperate father, and abandoned by respondents, who placed him in a dangerous predicament and who knew or learned what was going on, and yet did essentially nothing except, as the Court revealingly observes, \textit{ante} at 193, `dutifully recorded these incidents in [their] files.'''
\end{quote}
\noindent --- Justice Blackmun, dissenting, \href{https://supreme.justia.com/cases/federal/us/489/189/}{489 U.S. at 213}; \href{https://www.law.cornell.edu/supremecourt/text/489/189}{Cornell LII}.


\end{document}
